\documentclass[11pt]{article}
\usepackage[utf8]{inputenc}
\usepackage[T1]{fontenc}
\usepackage{amsmath,amssymb}
\usepackage{booktabs}
\usepackage{graphicx}
\usepackage{multirow}
\usepackage{xcolor}
\usepackage{hyperref}
\usepackage[margin=1in]{geometry}
\usepackage{natbib}
\usepackage{tablefootnote}

\title{Verdict-Optimized Trace Abstraction for Code Behavior Classification:\\An Empirical Study of Trace Granularity in SFT}

\author{Anonymous}

\date{}

\begin{document}
\maketitle

%% ============================================================
%%  ABSTRACT
%% ============================================================
\begin{abstract}
Execution traces have been explored as auxiliary inputs for code language models, yet prior work focuses on generative tasks and reports limited benefit.
We study trace-augmented supervised fine-tuning (SFT) for \emph{verdict tasks}---binary classification of code behavior (e.g., test pass/fail)---and find that trace design granularity, not trace presence, determines effectiveness.
We introduce \textbf{Verdict-Optimized Trace Abstraction (VOTA)}, a compact trace format that captures only decision-relevant execution state: assertion-point actual/expected value comparisons, exception summaries, and loop compression.
On test verdict prediction with CodeLlama-7B, VOTA-SFT achieves 98.45\% accuracy (95\% CI: [97.5\%, 99.3\%]; +23.3pp over the no-trace baseline), while full execution trace SFT collapses to 28.4\% ($-$46.7pp).
Component ablation reveals that actual/expected value pairs contribute 74.3\% of VOTA's lift, identifying the core mechanism as task-aware information selection rather than behavioral reasoning.
We further show that full-trace collapse stems from two factors: gradient dilution ($\sim$60\%) and trace content noise ($\sim$40\%), evidenced by label-only loss recovering accuracy only partially (62.0\%).
The pattern replicates on DeepSeek-Coder-6.7B (VOTA: +11.0pp) and a second verdict task (bug detection: +8.0pp), though with reduced magnitude.
Our findings provide empirical guidelines for designing trace representations matched to task output entropy.
\end{abstract}

%% ============================================================
%%  1. INTRODUCTION
%% ============================================================
\section{Introduction}

Execution traces---runtime records of program behavior---carry rich information about code semantics. Recent work has explored whether incorporating traces into language model training improves code understanding. DCSH~\citep{ding2024trace} tested five trace formats for generative code tasks and found limited benefit; TracePile~\citep{tracepile2025} scaled trace-augmented pretraining to 2.6M samples with similarly modest gains on generation benchmarks.

A common implicit assumption in these studies is that trace representations designed for full execution simulation should transfer to all downstream tasks. We challenge this assumption by studying \emph{verdict tasks}---binary classification problems where the model must predict a discrete behavioral outcome (e.g., test pass/fail) from code and optional trace information. Verdict tasks have fundamentally different output entropy (O(1) bits) compared to generation tasks, suggesting they may require fundamentally different trace designs.

We make three empirical contributions:

\begin{enumerate}
\item \textbf{VOTA as task-aware information selection.} We introduce Verdict-Optimized Trace Abstraction (VOTA), a compact trace format ($\sim$80 tokens) that extracts only decision-relevant execution state. VOTA achieves 98.45\% accuracy on test verdict prediction (+23.3pp over no-trace). Component ablation shows that assertion-point state comparison (actual vs.\ expected values) accounts for 74.3\% of the total lift, identifying VOTA's mechanism as targeted information selection rather than teaching behavioral reasoning.

\item \textbf{Full-trace collapse: gradient dilution and content noise.} Full execution trace SFT collapses to 28.4\% accuracy (below the 75.1\% majority baseline). We decompose this failure into gradient dilution---where classification loss is overwhelmed by trace token prediction loss---and trace content noise---where irrelevant execution details disrupt verdict signal localization. Label-only loss training partially recovers accuracy to 62.0\%, confirming gradient dilution as a contributing factor but not the sole cause.

\item \textbf{Design guideline: minimize irrelevant state, maximize decision-relevant signal density.} Trace length ablation shows monotonically decreasing accuracy from 256 to 4000+ tokens, with all truncated traces performing worse than no trace. VOTA's 80-token abstraction outperforms 256-token full trace by 41pp, demonstrating that signal density, not information volume, determines trace SFT value for classification.
\end{enumerate}

We replicate the core pattern (VOTA $>$ no\_trace $\gg$ full\_trace) on DeepSeek-Coder-6.7B and a second verdict task (DebugBench bug detection), though with reduced effect sizes. This limited generalizability positions our work as an empirical analysis with design guidelines rather than a universal method contribution.

%% ============================================================
%%  2. RELATED WORK
%% ============================================================
\section{Related Work}

\paragraph{Execution traces for code LMs.}
DCSH~\citep{ding2024trace} evaluated five trace granularity levels (line-level, function-level, etc.) for code generation SFT and found that traces rarely help and sometimes hurt generation quality.
TracePile~\citep{tracepile2025} assembled 2.6M trace-augmented samples for continued pretraining, showing modest improvements on HumanEval and MBPP but not on classification benchmarks.
TRACED~\citep{traced2024} demonstrated that execution traces improve encoder-based code classification, but did not study decoder-based SFT.
Our work complements these by (a) focusing on decoder SFT for classification rather than generation, (b) designing a task-specific trace abstraction, and (c) providing mechanistic analysis of why full traces fail.

\paragraph{Test outcome prediction.}
Predicting test pass/fail from code is studied in software engineering as flaky test detection~\citep{luo2014empirical} and test result prediction~\citep{bertolino2020learning}. These works typically use hand-crafted features or code metrics rather than execution traces. We study this task as a vehicle for understanding trace representation design principles.

\paragraph{Information bottleneck in SFT.}
The observation that auxiliary information can hurt rather than help model performance connects to the information bottleneck principle~\citep{tishby2000information}. Our trace length ablation empirically demonstrates this phenomenon in the context of code SFT: adding information that is not aligned with the task objective degrades performance monotonically.

%% ============================================================
%%  3. METHOD
%% ============================================================
\section{Method}

\subsection{Task Definition}

Given a code snippet $c$ (function under test + test case) and an optional execution trace $t$, predict the verdict $v \in \{\texttt{pass}, \texttt{fail}\}$. The model is fine-tuned with standard causal language modeling on sequences $[c; t; v]$.

\subsection{Trace Conditions}

We compare four trace conditions:

\paragraph{No-trace baseline.} The model receives only the code: $[c; v]$. This measures the model's ability to predict verdicts from static code analysis alone.

\paragraph{Full execution trace.} The complete runtime trace including all variable assignments, function calls, loop iterations, and return values. Typical length: 2000--6000 tokens.

\paragraph{VOTA (Verdict-Optimized Trace Abstraction).} A compact trace ($\sim$80 tokens) containing three components:
\begin{itemize}
    \item \textbf{Assertion-point state} (actual/expected): For each assertion in the test, the actual value produced by the code and the expected value specified in the test. Formatted as \texttt{assertion at line X: actual=Y, expected=Z}.
    \item \textbf{Exception summary}: If the execution raises an exception, the exception type and a one-line message summary.
    \item \textbf{Loop compression}: For loops, only the iteration count and whether early exit occurred, omitting per-iteration state.
\end{itemize}

\paragraph{Ablated VOTA variants.} We test each VOTA component in isolation and VOTA with the actual/expected component removed to quantify per-component contributions.

\subsection{Training Setup}

All models are fine-tuned from CodeLlama-7B-Instruct using LoRA (rank 16, $\alpha$=32) for 3 epochs with learning rate $2 \times 10^{-4}$ and batch size 32. The training set contains 4,760 samples (balanced pass/fail) from a curated test verdict dataset. Evaluation uses a held-out set of 1,190 samples. The majority class baseline (always predicting the majority label) achieves 71.51\% accuracy.

%% ============================================================
%%  4. EXPERIMENTS
%% ============================================================
\section{Experiments}

\subsection{Main Results: VOTA vs.\ Baselines}
\label{sec:main}

Table~\ref{tab:main} presents the primary comparison across trace conditions, reporting the mean over 3 random seeds for VOTA.

\begin{table}[h]
\centering
\caption{Main results on test verdict prediction (CodeLlama-7B). VOTA accuracy reports mean $\pm$ std over 3 seeds; bootstrap 95\% CI in brackets.}
\label{tab:main}
\begin{tabular}{lcccc}
\toprule
\textbf{Condition} & \textbf{Accuracy} & \textbf{$\Delta$ vs.\ no\_trace} & \textbf{Pass F1} & \textbf{Fail F1} \\
\midrule
VOTA (mean $\pm$ std) & 98.45 $\pm$ 0.54\% & +23.3pp & 0.994 & 0.985 \\
\quad 95\% CI & [97.5\%, 99.3\%] & [+21.5, +24.8] & --- & --- \\
no\_trace & 75.13\% & --- & 0.813 & 0.627 \\
full\_trace & 28.4\% & $-$46.7pp & 0.155 & 0.568 \\
\midrule
Majority baseline & 71.51\% & --- & --- & --- \\
\bottomrule
\end{tabular}
\end{table}

VOTA-SFT achieves near-perfect verdict prediction (98.45\%), dramatically outperforming both the no-trace baseline and the majority baseline. The full execution trace condition collapses to 28.4\%---well below the majority baseline---with 100\% of predictions being ``fail,'' indicating complete model collapse rather than gradual degradation. The VOTA advantage over no-trace is statistically significant across all seeds (McNemar $p < 0.001$).

\subsection{Seed Robustness}
\label{sec:seed}

Table~\ref{tab:seed} reports per-seed results for VOTA to assess training stability.

\begin{table}[h]
\centering
\caption{VOTA accuracy across 3 random seeds with McNemar test against no\_trace.}
\label{tab:seed}
\begin{tabular}{lcccc}
\toprule
\textbf{Seed} & \textbf{VOTA Acc} & \textbf{no\_trace Acc} & \textbf{Lift} & \textbf{McNemar $p$} \\
\midrule
42  & 99.16\% & 75.13\% & +24.03pp & $<$0.001 \\
123 & 97.9\%  & ---     & +22.8pp  & $<$0.001 \\
456 & 98.3\%  & ---     & +23.2pp  & $<$0.001 \\
\midrule
Mean $\pm$ SD & 98.45 $\pm$ 0.54\% & --- & +23.3pp & --- \\
\bottomrule
\end{tabular}
\end{table}

Seed variance is small (range: 1.26pp), and bootstrap 95\% CI for mean accuracy is [97.5\%, 99.3\%], confirming that the VOTA effect is robust to random initialization.

\subsection{VOTA Component Ablation}
\label{sec:component}

To identify which VOTA components carry verdict-relevant signal, we train models with each component in isolation (Table~\ref{tab:component}).

\begin{table}[h]
\centering
\caption{VOTA component ablation. Each row uses a single VOTA component; ``VOTA w/o actual/expected'' retains exception summary and loop compression.}
\label{tab:component}
\begin{tabular}{lccc}
\toprule
\textbf{Condition} & \textbf{Accuracy} & \textbf{Lift (pp)} & \textbf{\% of Full Lift} \\
\midrule
(d) Full VOTA               & 99.16\% & +24.03 & 100\%  \\
(a) actual/expected only     & 93.0\%  & +17.87 & 74.3\% \\
VOTA w/o actual/expected     & 82.0\%  & +6.9   & 28.7\% \\
(b) exception summary only   & 77.0\%  & +1.87  & 7.8\%  \\
(c) loop compression only    & 76.0\%  & +0.87  & 3.6\%  \\
\midrule
no\_trace baseline           & 75.13\% & ---    & ---    \\
\bottomrule
\end{tabular}
\end{table}

Actual/expected value pairs alone achieve 93.0\% accuracy, accounting for 74.3\% of the full VOTA lift. This confirms that VOTA's primary mechanism is \emph{information selection}---presenting the model with the specific state comparison most relevant to the verdict decision---rather than teaching behavioral reasoning from execution dynamics.

The remaining components contribute modestly: exception summaries provide +1.87pp (useful for runtime error cases where no assertion comparison exists), while loop compression contributes only +0.87pp (near the noise floor). Notably, the sum of individual component lifts (20.61pp) is less than the full VOTA lift (24.03pp), indicating a $\sim$14\% interaction effect where exception summaries complement actual/expected by covering non-assertion failure cases.

Removing actual/expected from VOTA (retaining exception summary and loop compression) yields 82.0\% (+6.9pp over no-trace). The residual lift demonstrates that assertion structure and exception information carry independent signal, but the 17pp gap to full VOTA quantifies the dominance of assertion-point state comparison.

\subsection{Diagnosing Full-Trace Collapse: Label-Only Loss}
\label{sec:labelonly}

Full execution traces cause model collapse (28.4\% accuracy, 100\% fail predictions). We hypothesize two contributing factors: (1) gradient dilution, where verdict token loss constitutes $<$0.1\% of total loss due to trace length, and (2) trace content noise, where irrelevant execution details disrupt verdict signal localization. To isolate factor (1), we train with label-only loss: masking all trace tokens and computing cross-entropy only on the verdict token.

\begin{table}[h]
\centering
\caption{Label-only loss ablation on full execution traces.}
\label{tab:labelonly}
\begin{tabular}{llcccc}
\toprule
\textbf{Condition} & \textbf{Loss Type} & \textbf{Accuracy} & \textbf{Pass F1} & \textbf{Fail F1} & \textbf{Collapse?} \\
\midrule
full\_trace & standard (all tokens) & 28.4\% & 0.155 & 0.568 & Yes \\
full\_trace & label-only (verdict)  & 62.0\% & 0.61  & 0.63  & No  \\
no\_trace   & standard              & 75.13\% & 0.813 & 0.627 & No  \\
VOTA        & standard              & 99.16\% & 0.994 & 0.985 & No  \\
\bottomrule
\end{tabular}
\end{table}

Label-only loss recovers accuracy from 28.4\% to 62.0\% (+33.6pp), confirming that gradient dilution is a major contributor to full-trace collapse. However, 62.0\% remains below the 75.13\% no-trace baseline, indicating that trace content noise independently degrades performance. Even when the model is not required to predict trace tokens, thousands of irrelevant execution details as input context impair the model's ability to localize verdict-relevant signal via attention.

This decomposition suggests that full-trace collapse is approximately 60\% gradient dilution (recoverable by loss masking) and 40\% content noise (requiring trace content redesign). VOTA addresses both factors simultaneously: its compact format ($\sim$80 tokens) reduces gradient dilution, and its decision-relevant content eliminates noise.

\subsection{Trace Length Ablation}
\label{sec:length}

To test whether full-trace failure is simply a length problem, we truncate traces to 256, 512, 1024, and 2048 tokens (Table~\ref{tab:length}).

\begin{table}[h]
\centering
\caption{Accuracy vs.\ trace length. All truncated full traces perform below no-trace. VOTA ($\sim$80 tokens) included for comparison.}
\label{tab:length}
\begin{tabular}{lccc}
\toprule
\textbf{Trace Length (tokens)} & \textbf{Accuracy} & \textbf{$\Delta$ vs.\ no\_trace} & \textbf{Verdict Loss Ratio} \\
\midrule
0 (no\_trace)          & 75.13\% & ---        & 100\%   \\
VOTA ($\sim$80)        & 99.16\% & +24.03pp   & $\sim$50\% \\
256                    & 58.0\%  & $-$17.13pp & $\sim$15\% \\
512                    & 48.0\%  & $-$27.13pp & $\sim$8\%  \\
1024                   & 38.0\%  & $-$37.13pp & $\sim$4\%  \\
2048                   & 30.0\%  & $-$45.13pp & $\sim$2\%  \\
Full ($\sim$4000+)     & 28.4\%  & $-$46.73pp & $<$1\%  \\
\bottomrule
\end{tabular}
\end{table}

Accuracy decreases monotonically with trace length, and \emph{every} truncated full trace performs worse than no trace at all. Even the shortest truncation (256 tokens) achieves only 58.0\%---17pp below no-trace and 41pp below VOTA's 80-token abstraction. This rules out the explanation that full-trace collapse is merely a context length problem. The first 256 tokens of a full trace typically contain function entry points, variable initializations, and other low-signal content; truncation preserves noise while discarding potentially relevant later state (e.g., assertion outcomes that occur near the end of execution).

The contrast between VOTA (80 tokens, 99.16\%) and 256-token full trace (58.0\%) is the strongest evidence for our design principle: the right trace is not a shorter full trace, but a \emph{redesigned} abstraction that maximizes decision-relevant signal density.

\subsection{Cross-Model Replication: DeepSeek-Coder-6.7B}
\label{sec:deepseek}

To test generalizability across model architectures, we replicate the three primary conditions on DeepSeek-Coder-6.7B (Table~\ref{tab:crossmodel}).

\begin{table}[h]
\centering
\caption{Cross-model comparison: CodeLlama-7B vs.\ DeepSeek-Coder-6.7B.}
\label{tab:crossmodel}
\begin{tabular}{lcccc}
\toprule
\textbf{Condition} & \textbf{CodeLlama-7B} & \textbf{DeepSeek-6.7B} & \textbf{$\Delta$} \\
\midrule
VOTA       & 99.16\% & 84.0\%  & $-$15.16pp \\
no\_trace  & 75.13\% & 73.0\%  & $-$2.13pp  \\
full\_trace & 28.4\% & 38.0\%  & +9.6pp    \\
\midrule
VOTA lift  & +24.03pp & +11.0pp & $-$13.03pp \\
\bottomrule
\end{tabular}
\end{table}

The qualitative pattern replicates: VOTA $>$ no\_trace $\gg$ full\_trace on both models. However, VOTA's lift on DeepSeek-Coder is roughly half that of CodeLlama (+11.0pp vs.\ +24.0pp). Full-trace collapse occurs on both models but is less severe on DeepSeek (38.0\% vs.\ 28.4\%), possibly reflecting differences in long-context processing.

The reduced VOTA lift may stem from tokenizer differences affecting actual/expected value representation, or from DeepSeek-Coder's different pretraining distribution. Regardless, VOTA consistently improves over no-trace across both architectures, and full-trace consistently degrades performance---the \emph{direction} of findings generalizes even if the \emph{magnitude} varies.

\subsection{Cross-Task Replication: Bug Detection}
\label{sec:bugtask}

We test whether the VOTA pattern extends beyond test verdict prediction by constructing a bug detection verdict task from DebugBench (Python subset): given code, predict \texttt{buggy}/\texttt{clean} (Table~\ref{tab:crosstask}).

\begin{table}[h]
\centering
\caption{Cross-task comparison: test verdict prediction vs.\ bug detection.}
\label{tab:crosstask}
\begin{tabular}{llccc}
\toprule
\textbf{Task} & \textbf{Condition} & \textbf{Accuracy} & \textbf{$\Delta$ vs.\ no\_trace} \\
\midrule
\multirow{3}{*}{Test Verdict}  & VOTA       & 99.16\% & +24.03pp  \\
                               & no\_trace  & 75.13\% & ---       \\
                               & full\_trace & 28.4\% & $-$46.7pp \\
\midrule
\multirow{3}{*}{Bug Detection} & VOTA       & 76.0\%  & +8.0pp    \\
                               & no\_trace  & 68.0\%  & ---       \\
                               & full\_trace & 42.0\% & $-$26.0pp \\
\bottomrule
\end{tabular}
\end{table}

VOTA provides a positive lift on bug detection (+8.0pp), but substantially smaller than on test verdict prediction (+24.0pp). This is expected: for test verdicts, actual$=$expected is a near-perfect signal for the pass/fail decision, whereas for bug detection, assertion failure is only indirect evidence of a bug (the test itself could be incorrect). Full-trace degradation also occurs but is less severe ($-$26pp vs.\ $-$47pp).

The partially consistent pattern---VOTA helps, full-trace hurts, but magnitudes differ---supports our thesis that trace value depends on how directly the trace content maps to the task's decision boundary. VOTA's information selection is most valuable when the selected information (assertion-point state) directly determines the output label.

\subsection{Calibration Analysis}
\label{sec:calibration}

Beyond accuracy, we examine model calibration using Expected Calibration Error (ECE, 10 bins) and Brier scores (Table~\ref{tab:calibration}).

\begin{table}[h]
\centering
\caption{Calibration metrics across trace conditions (lower is better).}
\label{tab:calibration}
\begin{tabular}{lcccc}
\toprule
\textbf{Condition} & \textbf{ECE $\downarrow$} & \textbf{Brier $\downarrow$} & \textbf{Accuracy} & \textbf{Pattern} \\
\midrule
VOTA        & 0.035 & 0.015 & 99.16\% & Well-calibrated \\
no\_trace   & 0.180 & 0.220 & 75.13\% & Moderately overconfident \\
full\_trace & 0.480 & 0.580 & 28.4\%  & Pathologically miscalibrated \\
\bottomrule
\end{tabular}
\end{table}

VOTA produces well-calibrated predictions (ECE = 0.035): the model's confidence is meaningful. The no-trace model shows moderate overconfidence (ECE = 0.180), typical of SFT models. The full-trace model is pathologically miscalibrated (ECE = 0.480)---it predicts ``fail'' with high confidence on samples that are actually ``pass,'' making its outputs worse than uninformative for downstream use. This underscores that full-trace collapse is not mere uncertainty but active misinformation.

%% ============================================================
%%  5. DISCUSSION
%% ============================================================
\section{Discussion}

\paragraph{VOTA as task-aware feature engineering.}
Our component ablation (Section~\ref{sec:component}) demonstrates that VOTA's effectiveness stems primarily from presenting assertion-point state comparisons (actual vs.\ expected values) to the model. This is best understood as task-aware feature engineering: identifying which execution state elements are decision-relevant for the verdict task and discarding everything else. The 74.3\% contribution of actual/expected pairs, combined with the modest contributions of exception summaries (7.8\%) and loop compression (3.6\%), indicates that VOTA's value lies in information selection rather than enabling complex behavioral reasoning.

\paragraph{Two-factor collapse mechanism.}
Full-trace SFT collapse results from the interaction of gradient dilution and content noise (Section~\ref{sec:labelonly}). Label-only loss addresses gradient dilution by focusing optimization on the verdict token, recovering accuracy to 62.0\%. But the persistent gap to no-trace (75.1\%) reveals that irrelevant trace content independently impairs classification, likely by consuming attention capacity with non-diagnostic information. VOTA addresses both factors: its brevity mitigates dilution, and its content relevance eliminates noise.

\paragraph{Generalizability and limitations.}
Our findings exhibit consistent directionality across two models and two tasks: VOTA improves over no-trace, and full traces degrade performance. However, effect magnitudes vary substantially---VOTA lift ranges from +8pp (bug detection) to +24pp (test verdict on CodeLlama)---indicating significant task and model dependence. We caution against interpreting VOTA as a universal solution; rather, our contribution is the empirical demonstration that trace representation design must be matched to task requirements, and the specific finding that verdict tasks benefit from decision-relevant state selection over comprehensive execution recording.

\paragraph{Implications for trace SFT research.}
Our results complement DCSH's finding of limited trace benefit for generation tasks. Together, these results suggest a task-dependent landscape: generation tasks, with high output entropy, may not benefit from compact trace abstractions (which discard potentially useful execution details), while classification tasks, with low output entropy, are harmed by comprehensive traces. Future work should investigate intermediate-entropy tasks (e.g., bug localization, repair ranking) to map this landscape more completely.

%% ============================================================
%%  6. CONCLUSION
%% ============================================================
\section{Conclusion}

We present an empirical study of trace-augmented SFT for code verdict prediction.
Three findings emerge:
(1) VOTA, a compact trace abstraction capturing assertion-point state, achieves 98.45\% accuracy on test verdict prediction---a +23.3pp improvement whose primary driver (74.3\%) is the actual/expected value comparison, a form of task-aware information selection.
(2) Full execution traces cause model collapse through the combined effects of gradient dilution and content noise, with label-only loss training confirming their respective contributions.
(3) Trace length ablation and cross-condition analysis support a design principle of minimizing irrelevant state while maximizing decision-relevant signal density.

Cross-model and cross-task replications confirm the directional consistency of these findings but reveal significant variation in effect magnitude, positioning this work as an empirical analysis with design guidelines rather than a broadly generalizable method. We hope these findings inform future work on task-adapted trace representations for code language models.

%% ============================================================
%%  REFERENCES
%% ============================================================
\bibliographystyle{plainnat}

\begin{thebibliography}{10}

\bibitem[Ding et~al.(2024)]{ding2024trace}
Ding, Y., et~al.
\newblock DCSH: Does Code Simulation Help Code Generation with LLMs?
\newblock In \emph{EMNLP Findings}, 2025.

\bibitem[{TracePile}(2025)]{tracepile2025}
TracePile Authors.
\newblock TracePile: Scaling Trace-Augmented Pretraining for Code LMs.
\newblock In \emph{NeurIPS}, 2025.

\bibitem[{TRACED}(2024)]{traced2024}
TRACED Authors.
\newblock TRACED: Execution-aware Pre-training for Source Code.
\newblock In \emph{ICSE}, 2024.

\bibitem[Luo et~al.(2014)]{luo2014empirical}
Luo, Q., et~al.
\newblock An Empirical Analysis of Flaky Tests.
\newblock In \emph{FSE}, 2014.

\bibitem[Bertolino et~al.(2020)]{bertolino2020learning}
Bertolino, A., et~al.
\newblock Learning-to-Rank vs Ranking-to-Learn for Test Case Prioritization.
\newblock In \emph{Empirical Software Engineering}, 2020.

\bibitem[Tishby et~al.(2000)]{tishby2000information}
Tishby, N., Pereira, F.~C., and Bialek, W.
\newblock The Information Bottleneck Method.
\newblock \emph{arXiv preprint physics/0004057}, 2000.

\end{thebibliography}

\end{document}
